\section{Adding a Single Anti-Fundamental Scalar} \label{sec:fundscalar}
We first extend the GG and BY models with one scalar field in the anti-fundamental representation. Thus, we introduce the scalar field $\phi$ to theory, which changes our Lagrangian. We will denote the chiral fermionic field as $\Lambda_{ij}$, referring either to $\Psi_{ij}$ for the case of the antisymmetric two-index tensor of the GG model or $\chi_{ij}$ in the case of the two-index symmetric tensor of the BY model. 

These field will also have charges under the $U_1(1)$ and the $\U_2(1)$ groups. These can be found by ensuring that we can write a Yukawa term for the chiral fermion, given by
\begin{equation}
    \mathcal{L}_{\text{Yukawa}}= y^\beta\Lambda_{ij} \tilde\psi ^{i}_{\beta} \phi^j + \mathrm{h.c.}\,.\label{eq:YukawaL}
\end{equation}
For this term to be charge-neutral, the total charge under each gauge group must vanish. For the first charge, $Q_1$, we have:
\begin{align}
     Q_{1_\Lambda} + Q_{1_{\tilde{\psi}}} + Q_{1_\phi} &= 0  \\
     (N \mp 4) - (N \mp 2) + Q_{1_\phi} &= 0 \quad \implies \quad Q_{1_\phi} = \pm 2 \,,
\end{align}
where the upper (lower) sign refers to the GG (BY) model. We will use this convention henceforth. Similarly, for the second charge, $Q_2$:
\begin{align}
     Q_{2_\Lambda} + Q_{2_{\tilde{\psi}}} + Q_{2_\phi} &= 0 \notag \\
     2p - p + Q_{2_\phi} &= 0 \quad \implies \quad Q_{2_\phi} = -p \,.
\end{align}
This charge assignment prohibits Yukawa terms mixing the vector-like fermions, since they are not invariant under $\U_1(1)$ and $\U_2(1)$. We will therefore have the Yukawa coupling written in \eqref{eq:YukawaL}.

The Lagrangian will also include self-interaction of the scalars. Ensuring renormalisability and symmetry invariance, we can write terms like
\begin{equation}
    \mathcal{L}_{\text{quartic}}\sim 
    \lambda \phi^\dagger_i \phi^i \phi^\dagger_j \phi^j + (\phi_i^\dagger (T^A)^i_{\,j}\phi^j)(\phi_k^\dagger (T^A)^k_{\,l}\phi^l)\,.\label{eq:quaritcL}
\end{equation}
However, for a scalar in the (anti-)fundamental representation, the terms are not linearly independent. This can be proved by considering that the generators $T^A$ satisfy the Fierz identity (see Appendix \ref{app:Fierz} for the derivation)
\begin{align}
    T_i ^{\,A\,j}T_k ^{\,A\,l}=\frac{1}{2}\left(\delta_i^{\,l}\delta_k^{\,j}-\frac{1}{N}\delta_i^{\,j}\delta_k^{\,l}\right) \,.
\end{align}
Applying this to the second term of \eqref{eq:quaritcL} yields
\begin{align}
    \frac{1}{2} (\phi_i^\dagger \phi^j)(\phi_k^\dagger \phi^l)\left(\delta_i^{\,l}\delta_k^{\,j}-\frac{1}{N}\delta_i^{\,j}\delta_k^{\,l}\right)=\frac{1}{2}\left[\phi^{\dagger}_i \phi^j \phi^{\dagger}_j \phi^i-\frac{1}{N}\phi^{\dagger}_i \phi^i \phi^{\dagger}_k \phi^k\right]=\frac{1}{2}\left(1-\frac{1}{N}\right) (\phi^{\dagger}_i \phi^i)^2\,,
\end{align}
Thus, the two quartic terms in the Lagrangian in Eq.~\eqref{eq:quaritcL} are not linearly independent, resulting in only one quartic term. 

Having collected all the pieces we can write the full Lagrangian as
\begin{align}
    \mathcal{L}= &-\frac{1}{4} F_{\mu\nu}^a F^{a\mu\nu}+ i \bar{\Lambda}^{ij} \slashed{D} \Lambda_{ij} 
    + i \bar{\psi}^i_\alpha \slashed{D} \psi_i^\alpha 
    + i \bar{\tilde{\psi}}_i^\beta \slashed{D} \tilde{\psi}^i_\beta\\\nonumber& +(D_\mu \phi)_i^\dagger (D^\mu\phi)^i+(y^\beta\Lambda_{ij} \tilde\psi ^{i}_{\beta} \phi^j + \mathrm{h.c.})+\lambda (\phi^\dagger_i \phi^i)^2 
    \,,
\end{align}
An overview of the model and the transformation of the fields can be seen in Table~\ref{tab:fund}.
\begin{table}[ht]
    \centering
    \renewcommand{\arraystretch}{1.5} 
    \begin{tabular}{|c|c|c|c|c|c|}
    \hline
    Fields & $\mathrm{SU}(N)$ & $\mathrm{SU}(N \mp 4 + p)$ & $\mathrm{SU}(p)$ & $\mathrm{U}_1(1)$ & $\mathrm{U}_2(1)$ \\
    \hline
    $\Lambda$ & $\ytableausetup{boxsize=0.25cm, centertableaux}\begin{ytableau}{}\\{}\end{ytableau} \;\, / \;\, \begin{ytableau}{}&{}\end{ytableau}$ & 1 & 1 & $N \mp 4$ & $2p$ \\

    $\tilde\psi$ & $\overline{\begin{ytableau}{}\end{ytableau}\rule{0pt}{7.5pt}}$ & $\overline{\begin{ytableau}{}\end{ytableau}\rule{0pt}{7.5pt}}$ & 1 & $-(N \mp 2)$ & $-p$ \\

    $\psi$ & ${\begin{ytableau}{}\end{ytableau}}$ & 1 & ${\begin{ytableau}{}\end{ytableau}}$ & $N \mp 2$ & $-(N - p)$ \\
    \hline
    $\phi$ & $\overline{\begin{ytableau}{}\end{ytableau}\rule{0pt}{7.5pt}}$ & 1 & 1 & $\pm 2$ & $-p$ \\
    \hline
    \end{tabular}
    \caption{Field content of the GG (upper signs, antisymmetric $\Lambda$) and BY (lower signs, symmetric $\Lambda$) models with an anti-fundamental scalar $\phi$.  The table shows how the fields transform under the gauge and global symmetry groups, and their charges under the $\mathrm{U}_1(1)$ and $\mathrm{U}_2(1)$ symmetries. }
    \label{tab:fund} 
\end{table}

It should be noted that the Yukawa term explicitly breaks the global symmetry $\mathrm{SU}(N\mp 4+p)$. The Yukawa couplings are a vector $y^\beta$ in the flavour space and, without loss of generality, it can be rotated using a $\mathrm{SU}(N\mp4+p)$ transformation so that only its first component is non-zero.\footnote{The flavour symmetry group is therefore broken following $\mathrm{SU}(N\mp4+p) \to \mathrm{SU}(N\mp 4+p-1)$, since the rotation of the Yukawa vector to one non-zero, is the same as selecting a specific component. Thus, only the remaining perpendicular components leaves the term invariant. } Hence, physically there is a single coupling constant, which we define as $y$, which greatly simplifies the computation of the beta functions since $y_\beta y^\beta=y^2$.

We begin with the RG equation for the gauge coupling. The only change from the computation done in the previous section is the inclusion of the anti-fundamental scalar, which contributes with $T(S) = 1/2$. Using \eqref{eq:betagauge}, we obtain 
\begin{equation}
\Delta b_0^{\text{scalar}} =  \frac{2}{3} T(S) =  \frac{1}{3}\,,\label{eq:scaleff}
\end{equation}
Thus, the beta function of the gauge coupling is given by
\begin{equation}
    \beta_g=-\frac{\alpha_g^2}{3}\left(18N-4p\pm{12}-1\right)\,.\label{eq:betagwithp}
\end{equation}
 The $(-1)$ inside the bracket corresponds to the scalar contribution.

The remaining beta functions are computed using the \texttt{RGBeta} library in Mathematica \cite{Thomsen:2021ncy}. Using similar rescaled couplings as for the gauge coupling $\alpha_y=y^2/(4\pi)^2$, $\alpha_\lambda=\lambda/(4\pi)^2$, the beta functions are given by
\begin{align}
    \beta_y&= \frac{\alpha_y}{2}\left[(3 N+2\mp3)\alpha_y+6 \alpha_g\left(-3 N\pm2+\frac 5N\right) \right]\,,\label{eq:betaywithp}
   \\
   \beta_\lambda&=\alpha_\lambda\left[2(N\mp1)\alpha_y-\frac{6 \alpha_g 
   \left(N^2-1\right)}{N}+4 \alpha_\lambda (N+4)\right]\nonumber\\&\;\;\;\;\;\;-\frac{1}{2} (N+1\mp 2 )
  \alpha_y^2+\frac{3 \alpha_g^2 \left(N^3+N^2-4 N+2\right)}{4 N^2}\,. \label{eq:betalamwithp}
\end{align}
\begin{figure}[b!]
    \centering
    \begin{subfigure}[b]{0.49\textwidth}
        \centering
        \includegraphics[width=\textwidth]{Figures/CAFGG_fund_Ng1.pdf}
        \caption{GG model} 
        \label{fig:cafgg_fund}
    \end{subfigure}
    \hspace{0.1cm}
    \begin{subfigure}[b]{0.49\textwidth}
        \centering
        \includegraphics[width=\textwidth]{Figures/CAFBY_fund_Ng1.pdf}
        \caption{BY model} 
        \label{fig:cafby_fund}
    \end{subfigure}
    \caption{CAF regions in parameter space $(N,p)$ for (a) the GG model and (b) the BY model, with one scalar in the anti-fundamental representation. The dark blue area represents the fixed-flow CAF region, bounded above by $b_0=0$ condition (dashed black line) and for the GG model the $b_0-d_2'+\sqrt{k'}=0$ below (dotted black line). The light blue region highlights the CAF region off fixed flow bounded by the $b_0=0$ condition above and the $k=0$ (dash-dotted black line) condition below. The striped region indicates that the model exhibits CAF both on and off fixed flow.  The boundaries are given by the $p_{\max}$ and $p_{\min}$ functions  Eqs. \eqref{eq:pmaxfundoff}, \eqref{eq:pminfundoff}, and  \eqref{eq:pminfundon}.}
    \label{fig:cafby}
\end{figure}
The limits of the CAF regions can be found by considering the boundary of the three conditions $b_0=0$, $k=0$, and $b_0-d_2'+\sqrt{k'}=0$ (see Eq.s \eqref{eq:CAF} and \eqref{eq:CAFfixed}). In the off-fixed-flow regime, the two former conditions give us 
\begin{align}
    p_{\text{max}}(N)&=\frac{18N\pm12-1}{4}\label{eq:pmaxfundoff}\,,\\
   p_{\text{min}}^\text{off}(N) &= -\frac{1}{4}  \pm 3  + \frac{9}{2N}  + \frac{3}{2N} \sqrt{ 3 (N+4)(N^3 + N^2 - 4N + 2) }\,.\label{eq:pminfundoff}
\end{align}
There is also an alternative branch of the $p_{\text{min}}^\text{off}(N)$ with a minus in front of the final term, which does not border the CAF region because the change results in a negative $p \,\forall\, N$. 

The boundary functions on fixed flow has the same upper boundary as in the off-fixed-flow regime (see Eq.~\eqref{eq:pmaxfundoff}). The lower boundary found from the $b_0-d_2'+\sqrt{k'}=0$ condition (see Eq. \eqref{eq:CAFfixed}) is given by
\begin{align}
    p_{\text{min}}^\text{on}(N)&= \frac{-720 + 4N(19 + 168N + 2N^2 - 9N^3)}{16N(N - 1)(N + 4)}\nonumber\\&\quad\;+\frac{ 3\sqrt{6}\sqrt{(1 - 3N)^2(N - 1)^2(N + 4)^2(N^2 + 2N - 2)}}{16N(N - 1)(N + 4)}\label{eq:pminfundon}\,,
\end{align}
where $p_{\text{min}}^\text{on}(N)$ only bounds the GG model, as can be seen from Figure~\ref{fig:cafgg_fund}.

We have now obtained the necessary beta functions and can continue the analysis of the models by performing a CAF analysis using the conditions mentioned in Eq.~\eqref{eq:CAF} and Eq.~\eqref{eq:CAFfixed}. We deploy these conditions on the beta functions Eqs. \eqref{eq:betagwithp}, \eqref{eq:betaywithp}, and \eqref{eq:betalamwithp}. It should be noted that the CAF regions found in this analysis \textit{may} correspond to fully asymptotically free theories. Whether these solutions correspond to CAF regions also depends on initial conditions of the couplings at a given RG scale (see Figure~\ref{fig:flowplottheory}). These conditions are provided in Eqs. \eqref{eq:Yukawacondoff}, \eqref{eq:ygfixed}, \eqref{eq:CAFalpha} and \eqref{eq:CAFalphafixed}, only when these are satisfied does one realise a genuine CAF regime. 

In Figure~\ref{fig:cafby}, the results of the CAF analysis are shown in the parameter space $(N,p)$. The light blue region in both figures represents the region where the coefficients of the beta functions satisfy the CAF conditions, when the Yukawa and gauge couplings are not on fixed flow. In other words, terms proportional to the Yukawa coupling have only a subleading effect in the beta function of the quartic couplings in this region.

Another feature of Figure~\ref{fig:cafby} is the dark blue region, which extends all the way up to the dashed black line. This is the region where the CAF conditions are satisfied for the Yukawa and gauge coupling on fixed flow. Thus, in this region, neither of the couplings is subleading. As mentioned in Section \ref{chap:CAF}, on fixed flow, the coefficients $d_2'$ and $d_4'$ can change sign. This is exactly what we observe in these models, effectively rendering the CAF conditions coming from the quartic coupling trivially satisfied. Thus, the fixed-flow region closely resembles that of the models without a scalar. 

To better understand these conditions and these regions, flow diagrams have been made both within and outside of each region. These can be found in Appendix \ref{app:flow}. The analysis gives the same result as the CAF analysis and clearly illustrates how the values of the couplings at the RG scale are restricted. 

\subsection{Inclusion of Chiral Families }\label{sec:fundscalarNg}
To further explore the GG  and BY models with a scalar in the fundamental representation, we now extend the analysis to include multiple chiral families. The transformation rules of the models will be similar to those before, as shown in Table~\ref{tab:fund_Ng}. The global symmetry group due to the multiplicity of anti-fundamental fermions $\mathrm{SU}(N\mp4+p)$ will change to $\mathrm{SU}(N_g(N\mp4)+p)$ ensuring anomaly cancellation (see Eq. \eqref{eq:anomSU}). Furthermore, another group, $\mathrm{SU}(N_g)$, is added, under which the chiral fermionic field transforms, effectively adding chiral families. This will also change the terms involving the (anti-)symmetric two-index fermionic field in the Lagrangian, by introducing a new flavour index $\gamma=1,\dots,N_g$, where $N_g$ is the number of chiral families. Putting this together, we find the following Lagrangian
\begin{align}
    \mathcal{L}&=  -\frac{1}{4} F_{\mu\nu}^a F^{a\mu\nu}+ i \bar \Lambda^{ij\gamma}\slashed{D}\Lambda_{ij\gamma}
    + i \bar{\psi}^i_\alpha \slashed{D} \psi_i^\alpha 
    + i \bar{\tilde{\psi}}_i^\beta \slashed{D} \tilde{\psi}^i_\beta +(D_\mu \phi)_i^\dagger (D^\mu\phi)^i\nonumber\\&\,\quad+(y^{\beta\gamma}\Lambda_{ij\gamma} \tilde\psi ^{i}_{\beta} \phi^j + \mathrm{h.c.})+\lambda (\phi^\dagger_i \phi^i)^2 
    \,,
\end{align}
where the Yukawa coupling is a $N_g\times (N_g(N\mp4)+p)$ matrix in flavour space.

We begin by considering the gauge beta function, which is given by
\begin{equation}
    \beta_g=-\frac{22N-1\pm 12 N_g-4N_gN -4p}{3}\alpha_g^2\,,
\end{equation}
where the beta function now depends on the number of chiral families $N_g$, together with $p$ and $N$. Hence, we can derive an upper limit on $p$ as a function of $N$ and $N_g$, namely
\begin{equation} \label{eq:pmaxfundNg}
    p_{\max} (N,N_g) = \frac{22N-1\pm12N_g-4N_g N}{4}\,.
\end{equation}
When $p_{\max} < 0$, no CAF region can exist as the gauge coupling is not asymptotically free, giving an upper bound on the number of chiral generations
\begin{equation}
    N_{g} \leq \frac{22N-1}{4(N\mp 3)}\,. 
\end{equation}
\begin{table}[ht]
    \centering
    \renewcommand{\arraystretch}{1.5} 
    \begin{tabular}{|c|c|c|c|c|c|c|}
    \hline
    Fields & $\mathrm{SU}(N)$ & $\mathrm{SU}(N_g(N \mp 4) + p)$ & $\mathrm{SU}(p)$ & $\mathrm{SU}(N_g)$ & $\mathrm{U}_1(1)$ & $\mathrm{U}_2(1)$ \\
    \hline
    $\Lambda$ & $\ytableausetup{boxsize=0.25cm, centertableaux}\begin{ytableau}{}\\{}\end{ytableau} \;\, / \;\, \begin{ytableau}{}&{}\end{ytableau}$ & $1$ & $1$ &$\overline{\begin{ytableau}{}\end{ytableau}\rule{0pt}{7.5pt}}$ & $N \mp 4$ & $2p$ \\

    $\tilde\psi$ & $\overline{\begin{ytableau}{}\end{ytableau}\rule{0pt}{7.5pt}}$ & $\overline{\begin{ytableau}{}\end{ytableau}\rule{0pt}{7.5pt}}$ & 1&1 & $-(N \mp 2)$ & $-p$ \\

    $\psi$ & ${\begin{ytableau}{}\end{ytableau}}$ & 1 & ${\begin{ytableau}{}\end{ytableau}}$ & 1&$N \mp 2$ & $-(NN_g - p)$ \\
    \hline
    $\phi$ & $\overline{\begin{ytableau}{}\end{ytableau}\rule{0pt}{7.5pt}}$ & 1 & 1 &1& $\pm 2$ & $-p$ \\
    \hline
    \end{tabular}
    \caption{Field content of the GG (upper signs, antisymmetric $\Lambda$) and BY (lower signs, symmetric $\Lambda$) models with $N_g$ chiral families and one anti-fundamental scalar $\phi$.  The table shows how the fields transform under the gauge and global symmetry groups, and their charges under the $\mathrm{U}_1(1)$ and $\mathrm{U}_2(1)$ symmetries. }
    \label{tab:fund_Ng} 
\end{table}
\begin{equation}
    \beta_{y}={\alpha_y}
    \left(3\alpha_g\left(\frac{5}{N}\pm2-3N\right)
    +\alpha_y\left(\mp N_g+N_gN+\frac{2\mp1+N}{2}\right)\right)\,,
\end{equation} 
where $N_g$ comes from the trace over Yukawa couplings in the beta function.\footnote{The trace factor is $N_g$ since $N_g \le N_gN\mp N_g4+p$ for all relevant  $N$ and $p$. } Using the same assumptions, the beta function of the quartic coupling is given by
\begin{align}
    \beta_{\lambda}=&\alpha_\lambda\left[2N_g(N\mp1)\alpha_y+\left(\frac{6}{N}-6N\right)\alpha_g+4 \alpha_\lambda (N+4)\right]-\frac{N_g}{2} (N+1\mp2)
  \alpha_y^2\nonumber\\&\,+\frac{3 \alpha_g^2 \left(N^3+N^2-4 N+2\right)}{4 N^2}\,.
\end{align}
These beta functions now include a dependency on the number of chiral families, effectively changing the CAF regions from the case of one chiral family in Section \ref{sec:fundscalar}.

We can again describe the boundaries of the CAF regions using maximal and minimum $p$ functions. The $p_{\max}$ function is given by the boundary of asymptotic freedom of the gauge coupling (see Eq.~\eqref{eq:pmaxfundNg}). For off-fixed-flow regimes the minimum is found from the $k=0$ condition (Eq. \eqref{eq:CAF}), which yields
\begin{align}
    p_{\min}^\text{off}(N,N_g)&= -\frac{1}{4} + N + \frac{3\left(3 + \sqrt{3}\sqrt{8 - 14N + 5N^3 + N^4}\right)}{2N} \pm 3N_g - N N_g\,.\label{eq:pminfundNg}
\end{align}
In the case of fixed flow we have the same $p_{\max}(N,N_g)$ function, i.e.\ Eq.~\eqref{eq:pmaxfundNg}, however the lower boundary differs between the two models. The BY models CAF regions extend all the way down to $p=0$. However, for the GG model the lower boundary from the condition $b_0-d_2'+\sqrt{k'}=0$ is found to be
\begin{align}
    p_{\min}^\text{on}(N,N_g)& = \frac{1}{16} \left( 68 + \frac{180}{N} - 20N + 48N_g - 16NN_g + \frac{3\sqrt{6}\sqrt{\frac{-2 + N(2+N)}{N_g}}(1 + N + 2(-1+N)N_g)}{N} \right)\label{eq:pminfundNgon}
\end{align}
The CAF region is demonstrated for both models in Figures  \ref{fig:CAFGGNG9} and \ref{fig:CAFNG9} in the $(N,p)$ parameter space.

\begin{figure}[t!]
    \centering
    % --- Top Row ---
    \begin{subfigure}[b]{0.49\textwidth}
        \centering
        \includegraphics[width=\textwidth]{Figures/CAFGG_NG1.pdf}
        \caption{$N_g=1$} 
        \label{fig:gg_ng1}
    \end{subfigure}%
    \hspace{0.1cm}
    \begin{subfigure}[b]{0.49\textwidth}
        \centering
        \includegraphics[width=\textwidth]{Figures/CAFGG_NG3.pdf}
        \caption{$N_g=3$}
        \label{fig:gg_ng3}
    \end{subfigure}
    
    \vspace{0.5cm} 
    % --- Bottom Row ---
    \begin{subfigure}[b]{0.49\textwidth}
        \centering
        \includegraphics[width=\textwidth]{Figures/CAFGG_NG6.pdf}
        \caption{$N_g=5$} 
        \label{fig:gg_ng6}
    \end{subfigure}%
    \hspace{0.1cm}
    \begin{subfigure}[b]{0.49\textwidth}
        \centering
        \includegraphics[width=\textwidth]{Figures/CAFGG_NG9.pdf}
        \caption{$N_g=6$}
        \label{fig:gg_ng9}
    \end{subfigure}
    
    \caption{CAF regions in the ($N$, $p$) parameter space for the GG model with $N_g \in \{1, 3, 5, 6\}$ with one anti-fundamental scalar. The dark blue area indicates the region of CAF on the fixed flow, which is bounded above by the $b_0=0$ condition (dashed black line) and below by the $b_0-c_1>0$ condition (dotted black line). The light blue area highlights the region of CAF off the fixed flow, bounded strictly between $b_0=0$ (above) and $k=0$ (dash-dotted black line, below) where the striped region indicates both cases exhibit CAF. The theoretical boundary curves correspond to the $p_{\max}$ and $p_{\min}$ functions detailed in Eqs. \eqref{eq:pmaxfundNg}, \eqref{eq:pminfundNg}, and \eqref{eq:pminfundNgon}.}
\label{fig:CAFGGNG9}
\end{figure}

\begin{figure}[t!]
    \centering
    % --- Top Row ---
    \begin{subfigure}[b]{0.49\textwidth}
        \centering
        \includegraphics[width=\textwidth]{Figures/CAFplotBY_NG1.pdf}
        \caption{$N_g=1$} 
        \label{fig:by_ng1}
    \end{subfigure}% <-- Keep this % sign!
    \hspace{0.1cm}
    \begin{subfigure}[b]{0.49\textwidth}
        \centering
        \includegraphics[width=\textwidth]{Figures/CAFplotBY_NG2.pdf}
        \caption{$N_g=2$} 
        \label{fig:by_ng2}
    \end{subfigure}
    
    \vspace{0.5cm} 
    
    % --- Bottom Row ---
    \begin{subfigure}[b]{0.49\textwidth}
        \centering
        \includegraphics[width=\textwidth]{Figures/CAFplotBY_NG3.pdf}
        \caption{$N_g=3$} 
        \label{fig:by_ng3}
    \end{subfigure}% <-- Keep this % sign!
    \hspace{0.1cm}
    \begin{subfigure}[b]{0.49\textwidth}
        \centering
        \includegraphics[width=\textwidth]{Figures/CAFplotBY_NG4.pdf}
        \caption{$N_g=4$} 
        \label{fig:by_ng4}
    \end{subfigure}
    
    \caption{CAF regions in the ($N$, $p$) parameter space for the BY model for $N_g \in \{1, 2, 3, 4\}$ with one anti-fundamental scalar. The dark blue area indicates the region of CAF on the fixed flow, which is bounded above by the $b_0=0$ condition (dashed black line). The light blue area highlights the region of CAF off the fixed flow, bounded strictly between $b_0=0$ (above) and $k=0$ (dash-dotted black line, below) where the striped region indicates both cases exhibit CAF. The theoretical boundary curves correspond to the $p_{\max}$ and $p_{\min}$ functions detailed in Eqs. \eqref{eq:pmaxfundNg} and \eqref{eq:pminfundNg}}
    \label{fig:CAFNG9}
\end{figure}

As visualised in Figure~\ref{fig:CAFGGNG9} the CAF region becomes smaller as $N_g$ becomes larger, due to the positive contribution of the fermions to the $b_0$ coefficient. While the GG model is explicitly shown for $N_g=(1,3,5,6)$ it can accommodate CAF regimes with integer solutions up to $N_g={13}$. 

It is of interest to consider the case of $\mathrm{SU}(5)$ with three chiral families, which corresponds to a GUT with the Higgs multiplet scalar \cite{Georgi:1974yf} (see Section \ref{sec:SU5}). In this setup the fundamental scalar $\phi$ can break the electroweak symmetry. Our analysis shows that the GG model possesses CAF regions both on and off fixed flow. Specifically, when the Yukawa coupling is on fixed flow the smallest amount of vector-like fermions needed to realise CAF is $p=4$. However, when the Yukawa coupling is off fixed flow the requirement increases to $p=18$. This highlights the viability of the $\mathrm{SU}(5)$ GUT when supplied with a sufficient vector-like sector. 

The BY model is shown in Figure~\ref{fig:CAFNG9}. The figure displays four values of $N_g=(1,2,3,4), $ where for $N_g\ge 5$ the CAF region vanishes contrary to the GG model. In all the cases the theory exhibits CAF both on and off fixed flow. 











